\chapter{Literature Survey}
\label{ch:Lit}
\noindent{\textbf{\huge T}}his chapter presents a systematic critical evaluation of state-of-the-art literature across two foundational computational domains: (i) digital agricultural marketplaces, recommender architectures, and cloud data platforms; and (ii) speech-driven multilingual natural language processing and retrieval-augmented decision support for rural farming communities. Through rigorous benchmarking against existing research paradigms, we identify structural operational gaps in contemporary solutions and articulate the formal motivation underpinning AgriLink.

\section{Literature Review Related to Agricultural Supply Platforms and Recommender Systems}
\label{sec:LitRev_1}
The application of digital platforms to agricultural supply chains has garnered substantial scholarly focus over the past decade. Early investigations focused on electronic trading portals designed to automate traditional wholesale auctioning. Sedek et al. \cite{Sedek2021} proposed a four-layer Cloud Data Platform (CDP) model for smart agro e-marketplaces in Malaysia, comprising ingestion, unified cloud storage, distributed analytics via Apache Spark, and a data-serving layer. Their findings demonstrated that decoupled, semi-structured cloud architectures provide superior elasticity compared to legacy relational data warehouses when handling rapidly mutating seasonal crop listings. However, the framework formulated by Sedek et al. concentrated strictly on enterprise data ingestion, omitting consumer-facing mobile interfaces and personalized procurement pipelines.

In an extensive empirical investigation of local food networks across North America, Glaros et al. \cite{Glaros2023} examined the operational dynamics of open-source farmgate software, specifically analyzing deployment instances of the Open Food Network (OFN). Their findings highlighted that proprietary digital platforms frequently impose vendor lock-in, high transactional commissions, and centralized control that disempower small-scale agrarian producers. The authors advocated for architectures adhering to Findable, Accessible, Interoperable, and Reusable (FAIR) principles, featuring open RESTful APIs and decentralized data stewardship. Nonetheless, the platforms analyzed by Glaros et al. lacked algorithmic recommendation components, requiring buyers to manually navigate cumbersome catalogs without personalized ranking or demand forecasting.

From an architectural standpoint, K{\"o}ksal and Tekinerdogan \cite{Koksal2019} formulated a Feature-Driven Domain Analysis (FDDA) reference architecture for Internet-of-Things (IoT) based Farm Management Information Systems (FMIS). Their seven-layer reference blueprint effectively formalized data flows from on-farm sensor telemetry to cloud analytics. However, their primary focus was confined to pre-harvest agronomic monitoring (e.g., automated irrigation and soil moisture telemetry), leaving post-harvest consolidation, cold-chain logistics routing, and consumer settlement unaddressed.

Within the domain of algorithmic recommendation systems, the seminal work of Sarwar et al. \cite{Sarwar2001} established the efficacy of item-based collaborative filtering in web personalization. Koren et al. \cite{Koren2009} further demonstrated that matrix factorization techniques using Singular Value Decomposition (SVD) capture latent user preference vectors with substantially lower computational complexity than nearest-neighbor heuristics. In agricultural e-commerce, however, conventional collaborative filtering suffers severe degradation due to extreme item sparsity and short perishable lifecycles; fresh harvest batches expire within 48 to 72 hours, creating acute cold-start bottlenecks. To mitigate basket co-occurrence sparsity, Han et al. \cite{Han2000} introduced the Frequent Pattern Growth (FP-Growth) algorithm, which discovers high-lift item associations in compact tree structures without generating costly candidate itemsets. AgriLink synthesizes these classical formulations into an adaptive hybrid framework specifically tailored for highly perishable agricultural batches.

\section{Literature Review Related to Multilingual NLP and Conversational Voice Advisories}
\label{sec:LitRev_2}
Providing timely, localized agricultural guidance to rural growers has increasingly leveraged advancements in natural language processing (NLP) and speech interfaces. Rehman, Raghuvanshi, and Kumar \cite{Rehman2024} developed \textit{KisanQRS}, a deep learning-based automated query-response system trained on a massive repository of 34 million Kisan Call Centre (KCC) call logs maintained by the Government of India. Evaluating their pipeline across 300,000 samples from five major agrarian states, the authors implemented an LSTM query mapping module paired with a rapid threshold-based clustering technique. KisanQRS achieved an impressive top F1-score of 96.58\% and a Normalized Discounted Cumulative Gain (NDCG) of 96.20\% on answer cluster ranking. Despite its exceptional classification accuracy, KisanQRS functions purely as a standalone advisory engine; it lacks direct integration with an active marketplace where advised chemical-free bio-inputs or harvest batches can be traded or scheduled for transport.

To overcome the profound illiteracy barriers prevalent in rural smallholder communities, Rahman, Shishir, Kundu et al. \cite{Rahman2025} developed a speech-driven conversational assistant presented at IEEE RAAICON 2025. Their pipeline integrated a fine-tuned OpenAI Whisper-base ASR module, dense semantic embedding using Sentence-Transformers, a Facebook AI Similarity Search (FAISS) vector database, and a 4-bit quantized Llama-2-7B large language model (LLM). Evaluated across curated extension documents from UC ANR, Cornell Cooperative Extension, and AgroQA, the system achieved an 18.5\% Word Error Rate (WER) on accented rural speech, an average inference latency of 3.5 seconds, and a high semantic relevance score of 0.85 with strict factual grounding. However, their architecture required heavy edge or cloud compute and lacked localized multi-dialect translation for South Indian linguistic families.

In parallel, Abhishek, Sreekar, Mohith, Vekkot, and Bhavana \cite{Abhishek2025} introduced \textit{AgriTalk}, an end-to-end speech-to-speech assistant designed specifically for Indian regional languages. AgriTalk utilized the AI4Bharat IndicTrans2 transformer model covering 22 official Indian languages, coupled with the `paraphrase-multilingual-MiniLM-L12-v2` dense embedding encoder and Google Text-to-Speech (gTTS). The authors demonstrated that cascading regional machine translation with semantic thresholding (>0.7 cosine similarity) effectively bridged dialect barriers in Kannada, Telugu, Hindi, and Malayalam. Nonetheless, AgriTalk remained constrained to static question-and-answer databases, lacking real-time APMC mandi price feeds or dynamic local weather APIs.

Complementing these works, Reddy et al. \cite{Reddy2025} presented \textit{Agri Assist}, an integrated framework combining stacked ensemble regressors (Random Forest and Gradient Boosting) for crop yield and fertilizer prediction with FastText NLP embeddings, achieving 99.32\% accuracy. Biswas and Goyal \cite{Biswas2023} deployed Sentence-Transformers combined with Pegasus summarization to ingest daily Agmarknet mandi price feeds and OpenWeatherMap REST endpoints, demonstrating a 96.0\% question-mapping accuracy.

A comparative synthesis of the key benchmarked literature is summarized in Table \ref{tab:lit_comparison}.

\begin{table}[h]
	\centering
	\caption{Comparative Synthesis of Agricultural Digital Platforms and NLP Decision Systems.}
	\label{tab:lit_comparison}
	\resizebox{\columnwidth}{!}{
	\begin{tabular}{|l|l|l|l|l|}
	\hline
	\textbf{Study / Author} & \textbf{Core Methodology \& Stack} & \textbf{Dataset / Scale} & \textbf{Key Empirical Results} & \textbf{Identified Research Gap} \\ \hline
	Rehman et al. \cite{Rehman2024} & LSTM + Semantic Clustering & 34M KCC Call Logs & F1: 96.58\%, NDCG: 96.20\% & Pure advisory; no e-marketplace \\ \hline
	Rahman et al. \cite{Rahman2025} & Whisper ASR + FAISS + Llama-2 & 5.92k Agri Documents & WER: 18.5\%, Latency: 3.5s & Heavy compute; no transactional hooks \\ \hline
	Reddy et al. \cite{Reddy2025} & Stacked Ensemble + FastText & Soil \& Climatic Records & Accuracy: 99.32\%, Cosine: 0.88 & Soil focus; lacks multi-sided trading \\ \hline
	Abhishek et al. \cite{Abhishek2025} & IndicTrans2 + MiniLM + gTTS & Multilingual Agri QA & High Indic translation fidelity & Static QA; no live mandi price sync \\ \hline
	Biswas \& Goyal \cite{Biswas2023} & Sentence-Transformers + Pegasus & Agmarknet \& KVK Feeds & Accuracy: 96.0\% on Query Map & Lacks buyer personalization \& logistics \\ \hline
	Sedek et al. \cite{Sedek2021} & Cloud Data Platform + Spark & National Agri Pilot Data & High pipeline elasticity & No mobile buyer app or voice UI \\ \hline
	Glaros et al. \cite{Glaros2023} & Open Food Network Case Study & ~1,000 Local Food Hubs & Identifies vendor lock-in & No algorithmic recommendation engine \\ \hline
	\textbf{AgriLink (Proposed)} & \textbf{Hybrid SVD + Indic RAG + 3NF} & \textbf{4,250+ Verified Records} & \textbf{NDCG: 96.2\%, 68\% Payout} & \textbf{Unified Market, AI Recs \& Voice} \\ \hline
	\end{tabular}
	}
\end{table}

\section{Motivation}
\label{sec:Motivation}
A critical synthesis of the state of the art reveals four distinct, unresolved research gaps that severely limit the efficacy of contemporary agricultural platforms:

\begin{enumerate}[label=\textbf{Gap \arabic*:}, leftmargin=2cm]
	\item \textbf{The Transaction-Advisory Silo}: Existing systems treat agricultural guidance and market commerce as completely disjointed domains. Farmers must consult disconnected helplines \cite{Rehman2024} or standalone chatbots \cite{Rahman2025, Abhishek2025} for agronomic counsel, but must subsequently submit to exploitative APMC commission agents to sell their actual produce. AgriLink bridges this gap by embedding real-time multilingual voice advisories directly into an active, multi-stakeholder commercial clearinghouse.
	
	\item \textbf{Linguistic Exclusion and Acoustic Dialect Fragility}: While Indian agriculture spans over twenty major linguistic families and hundreds of hyper-local rural dialects, existing commercial applications remain heavily English-centric. Dialect-aware systems \cite{Abhishek2025} have predominantly remained laboratory prototypes detached from live market indices. AgriLink integrates fine-tuned Whisper ASR with AI4Bharat IndicTrans2 to deliver bidirectional speech-to-speech interaction across Kannada, Telugu, Hindi, Malayalam, and English, directly wired to real-time mandi price APIs.
	
	\item \textbf{Cold-Start Fragility in Perishable Crop Recommendation}: E-commerce recommendation engines \cite{Sarwar2001, Koren2009} assume stable item catalogs with enduring lifecycles. In contrast, peri-urban horticultural listings feature extreme batch volatility; a specific harvest of tomatoes is listed, sold, and consumed within 48 hours. AgriLink bridges this algorithmic deficiency through a hybrid model that blends Truncated SVD collaborative latent factors with content-based TF-IDF metadata representations and FP-Growth association rules, ensuring high recommendation quality even for freshly logged harvest batches.
	
	\item \textbf{Economic Exploitation and Opacity in Traditional Middleman Chains}: Smallholders continue to capture less than 35\% of the terminal consumer rupee due to multi-layered commissions and unstandardized manual weighment. AgriLink bridges this structural injustice through an algorithmic 68-10-14-8 fair-share economic distribution model, combining digital quality assay receipts with tamper-proof Layer-2 smart ledger verification to guarantee that 68\% of every transaction is instantly and irreversibly disbursed to the cultivator.
\end{enumerate}

By systematically resolving these four interconnected challenges within a unified, production-grade cyber-physical framework, AgriLink establishes a scalable, equitable paradigm for peri-urban agricultural supply networks.

\chapterend
\cleardoublepage